%-----------------------------------
\section{Trabalhos relacionados}
\label{sec_trabalhos_relacionados}

% ---------------------------------------------------------------
% ORIENTAÇÃO: apresente de três a seis trabalhos próximos ao seu,
% um parágrafo para cada (objetivo, método, principais resultados),
% e termine com um quadro comparativo que evidencie o diferencial
% do seu trabalho. Em revisões sistemáticas, esta seção é curta ou
% inexistente, pois a própria revisão cumpre esse papel.
% ---------------------------------------------------------------

Esta seção apresenta trabalhos que abordam problemas semelhantes ao investigado. Para cada trabalho, descreva em um parágrafo o objetivo, o método empregado e os principais resultados, sempre com a devida citação. Em seguida, compare os trabalhos segundo critérios relevantes para a sua pesquisa, como no Quadro~\ref{qua_comparativo}, e explique o que diferencia a sua proposta.

\begin{quadro}[htb]
	\centering
	\caption{Comparação entre os trabalhos relacionados e este trabalho (exemplo).}
	\label{qua_comparativo}
	\small
	\begin{tabular}{|l|c|c|c|c|}
		\hline
		\textbf{Trabalho} & \textbf{Critério 1} & \textbf{Critério 2} & \textbf{Critério 3} & \textbf{Critério 4} \\ \hline
		Trabalho A & Sim & Não & Sim & Não \\ \hline
		Trabalho B & Sim & Sim & Não & Não \\ \hline
		Trabalho C & Não & Sim & Sim & Não \\ \hline
		\textbf{Este trabalho} & \textbf{Sim} & \textbf{Sim} & \textbf{Sim} & \textbf{Sim} \\ \hline
	\end{tabular}
	\fonte{Elaborado pelo autor.}
\end{quadro}
